% !TeX root = main.tex
\chapter{有限域}

\section{域的定义}

\begin{definition}
    设\(R\)是一个环，若\(R\)中除零元以外的元素都具有乘法逆元，则称\(R\)是一个域。
    若\(R\)是一个域且\(R\)中元素的个数只有有限的\(p\) 个，则称\(R\)是一个有限域，
    记作\(\mathbb{F}_{p}\) 或\(\mathrm{GF}(p)\)。
\end{definition}

% 例子 Zp，证明Zp是域

\begin{theorem}[]
    域是整环
\end{theorem}

\begin{proof}
    反证法。若存在\(a,b\neq 0\)使得\(ab=0\)，则\(b=a^{-1}ab=a^{-1} \cdot 0=0\)，
    这与\(b\neq 0\)矛盾。
\end{proof}

\begin{theorem}
    有限整环一定是域
\end{theorem}

\begin{proof}
    设\(R\)是一个有限整环，任取\(a\neq 0\)，
    考虑映射\(f(x)=ax\)。由于\(R\)是一个整环，若\(x\neq y\),则\(ax\neq ay\)。故\(f\)是单射。
    又因为\(R\)是有限的，所以\(f\)也是满射。因此，存在一个元素\(b\in R\)，使得\(ab=f(b)=1\)。因此，
    \(a\)的逆元为\(b\)，\(R\)是一个域。
\end{proof}

但是无限整环不一定是域，例如\(\mathbb{Z}\)。

\section{域的特征}

\begin{definition}[域的特征]
    设\(F\)是一个域，若存在一个最小正整数\(p\)，使得\(p \cdot e = 0\)，则称\(F\)的特征为\(p\)。
    否则，称\(F\)的特征为\(0\)。
\end{definition}

\begin{theorem}[域的特征为素数]
    设\(F\)是一个域，若\(F\)的特征为\(p \)，则\(p\)是一个素数。
\end{theorem}

\begin{proof}
    反证法。设\(p=p_{1}p_{2}\)，其中\(p_1,p_2>1\)。则\(0=p \cdot e=(p_{1}p_{2}) \cdot
    e=(p_{1}e) \cdot (p_{2}e)\)，
    这与\(F\)是整环矛盾。
\end{proof}

\begin{theorem}[]
    设\(F\)是特征为\(p\)的域，则\(\forall a \in F\)，有\(p \cdot a = 0\)。

    进一步，有\(ma=0 \iff p \mid m\)
\end{theorem}

\begin{proof}
    设\(a\in F\)，则\(0=p \cdot e=(p \cdot e) \cdot a=p \cdot (e \cdot
    a)=p \cdot a\)。

    设\(m=kp+r\)，其中\(k,r\in \mathbb{Z}\)，且\(0 \leq r < p\)。则
    \(ma=(kp+r)a=k(p \cdot a)+ra=ra\)。因此，\(ma=0\)当且仅当\(r=0\)，即\(p \mid m\)。
\end{proof}

\begin{theorem}[]
    设\(F\)是特征为\(p\)的域，则\(\forall a,b \in F\)，有\[
        (a\pm b)^{p}=a^{p}\pm b^{p}
    \]
\end{theorem}

\begin{proof}
    设\(a,b \in F\)，则\[
        (a+b)^{p}=\sum_{k=0}^{p}\binom{p}{k} a^{k} b^{p-k}
    \]
    因为\(p\)是素数，所以当\(1 \leq k \leq p-1\)时，
    \(\binom{p}{k}=\frac{p!}{k!(p-k)!}\)是\(p\)的倍数。
    因此，\[
        (a+b)^{p}=a^{p}+b^{p}
    \]

    同理可得\((a-b)^{p}=a^{p}-(-1)^{p}b^{p}\)。
    \begin{itemize}
        \item 当\(p=2\)时，\(b^{p}= - b^{2} \)，\((a-b)^{p}=a^{p}-b^{p}\)。
        \item 当\(p\)为奇素数时，\((-1)^{p}=-1\)，\((a-b)^{p}=a^{p}-b^{p}\)。
    \end{itemize}
\end{proof}
% todo: 为什么组合数一定是整数
\begin{corollary}
    设\(F\)是特征为\(p\)的域，则\(\forall a ,b \in F\)，有\[
        (a\pm b)^{p^{n}}=a^{p^{n}}\pm b^{p^{n}}
    \]

    \(\forall a_{1},\dots ,a_{k} \in F\)，有\[
        (a_{1}+\cdots +a_{k})^{p}=a_{1}^{p}+\cdots +a_{k}^{p}
    \]
\end{corollary}

我们通过研究子域 \(\langle e \rangle\)来研究域的性质。

\begin{definition}[素域]
    设\(e\)是域\(F\)的乘法单位元，考虑：\[
        \Pi \coloneqq \left\{ 0, e, 2e ,\dots (p-1)e \right\}
    \]
    则\(\Pi\)是\(F\)的子域，称为\(F\)的素域。
\end{definition}

\begin{proof}
    \(\forall 0\leq k< p, 0\leq l<p, k\neq p\)，则：
    \begin{align*}
        ke-le & = (k-l)e \in \Pi \\
        (ke)(le)^{-1}&= (kl^{-1})e \in \Pi
    \end{align*}
    其中的减法与逆是在模\(p\)意义下进行的。
\end{proof}

素域是域的最小子域，同构于\(\mathbb{Z}_{p}\)。

\section{子域/域同构}

域同构与子域的概念和群/环基本类似，这里不加证明的给出一些定理。

\begin{theorem}[子域的判定]
    设\(F\)是一个域，\(S\subseteq F\)，则\(S\)是\(F\)的子域当且仅当：
    \begin{itemize}
        \item \(e \in S \)
        \item \(\forall a,b \in S\)，有\(a-b \in S\)和\(ab^{-1} \in
            S\)（\(b \neq 0\)）。
    \end{itemize}
\end{theorem}

\begin{theorem}[域同构]
    设\(F\)和\(F'\)是两个域，\(f:F \to F'\)是一个映射，则\(f\)是\(F\)到\(F'\)的同构当且仅当：
    \begin{itemize}
        \item \(f\)是双射。
        \item \(\forall a,b \in F\)，
            有\(f(a+b)=f(a)+f(b)\)和\(f(ab)=f(a)f(b)\)。
    \end{itemize}
\end{theorem}
同构的域特征相同。

theorem: a^{p^{n}} 是域上的自同构
\section{域上的多项式}

\begin{theorem}[域上的多项式环]
    \(\langle F[x], +, \cdot \rangle\)是带幺交换环。
\end{theorem}

% 分裂域，同构定理
% 多项式环的理想与商环

% 唯一因式分解怎么证明的？
% 多项式无重因式当且仅当其与其导数互素。
% 因式定理是余元定理的导出

% 只要能做带余除法，是不是理想就形如<elem>?
% F[x]/<p(x)> 原来就是和Z/nZ一样的感觉

% 有限域乘法部分一定是个循环群。从而可以定义本原元，并用欧拉计数公式计算个数
% 其中证明用了多项式，特别神奇的是域与他的多项式环看起来没关系，但多项式环上的定理却能限制域本身的结构。这是正确的
% 通过抽离有限域的元素，可以说明有限域本质上就是Zp的扩域。

% 有限域本质上是Zp通过扩张若干个元得到的。但是这若干个元不可说，也找不到。因此只能用同构的方式模去一个n次不可约多项式得到。
% 但这些多项式也不是随便来的，都是x^{p^n} - x的因式。（这本质上又是乘法循环群的拉格朗日定理

% 而且这个系统和特征也有关系，因为在加法与单位元构成的域上，F肯定是他的一个线性空间，因此维数是有限的。又因为域的特征为p，
% 所以这个线性空间的基底元素的个数一定是p的幂次。

% 莫比乌斯反演与偏序集

% Fq*是一个乘法循环群，这就导致a^q - a =0。然后域会有子域，域的大小一定是子域p 的某个次方，也就是 p^n，
% 这是用不断扩张元素做到的。这里书上不断的用 Fq[x] / <p(x)> 这种写法，我的看法就类似于说，p(x)的根不在这里面，
% 我知道Fq是靠底层的 Fp扩张若干个元素做到的，但是我又不能说出来，我就把他看作是p(x)的根。

% 域的特征是会传染的，底层域的特征是p，扩张域的特征也是p。因为扩张域是底层域的线性空间。

% 极小多项式是域外面的元素在域里面东西的体现。
% 在域里面，你可能感觉特别神秘，但是在域外面，一个东西的极小多项式本质上就是（x-α）的乘积，其中α是该元素在域外面的共轭。

% 本原元的好处就在于把所有的东西都变成了本原元的幂

% 在证明域的乘法群是循环群的时候，我们取域中阶最大的那个，然后证明它的阶一定是域的大小减一。本质就是取本原元，域的乘法群就是本原元生成的循环群。

% 有限阿贝尔群基本定理，$G$ 同构于若干个循环群的直积：
% $G \cong \mathbb{Z}_{d_1} \times \mathbb{Z}_{d_2} \times \dots \times \mathbb{Z}_{d_k}$

% 前缀码的解析是靠树实现的
% 比如说我想实现一个自己的str，或者long int，那我应该占用内存一部分并直接写入二进制数据。那这种需求，怎么在一般的编程语言中实现？

% 为什么三棱柱就能三线合一？

% 汉明球看起来像球，实际上也不是球体，因为看似距离是｜x-u｜《 n, 实际上是 x!= u. 球里面元素个数与x无关，
% 不会出现在边边角角的元素球体积就小的问题

% 每一列乘上某一个元是不是就是所有的换元型编码？
% 身份证号、手机号码都是编码
% 微分算子构成的多项式是ED吗
% 有限域子空间计数问题我总觉得和群有关系
